% GENERATED FILE. Edit canonical lessons, then run npm run generate:books.
% canonical-source-hash: fnv1a64:2733d56e824c2bf6
% canonical-lessons: MR-C31-naa-naa, MR-C31-naa-tag, MR-C31-dusra, MR-R31-pairs

\chapter{Words That Need Two Places}
\label{ch:mr-paired-joins}
\hlchaptermodality{31}

\begin{chapteropening}
\textbf{By the end of this chapter:} \emph{I can refuse two things at once with naa ... naa, check a statement with a final naa, and share a pair out with ek ... dusraa.}

Complete MR-R31-pairs and run chaha and dudh three ways: refuse both, share them out, and check a statement about tea.
\end{chapteropening}

\section[nā … nā …]{\mr{ना} … \mr{ना} … (nā … nā …) — neither, nor}

\label{lesson:MR-C31-naa-naa}



\begin{quote}
Retrieve the last two joining words by name: the one that offers a choice (\textbf{\mr{किंवा}}) and the one that takes something back (\textbf{\mr{पण}}). Say a sentence with each before you go on.
\end{quote}

\begin{grammarlens}[title={Grammar Lens: the word that comes in a pair}]
\begin{quote}
\textbf{\mr{ना} \mr{चहा} \mr{ना} \mr{दूध}.} — \emph{nā chahā nā dūdh.} — \textbf{Neither tea nor milk.}
\end{quote}

Every joining word so far went in ONE place, between the halves. This one goes in \textbf{two}, in FRONT of each half, and it must be said both times. Say it once and the sentence is unfinished.

\textbf{\mr{ना}} is the short form of \textbf{\mr{नाही}}, the \emph{no} you have had since the courtesy chapter, and it carries the same ancient \textbf{na-}. Doubling it is what makes it a joiner rather than a plain refusal: the first \textbf{\mr{ना}} sweeps the first thing away, the second sweeps the second, and the pair means \emph{both of them, gone}.

English does exactly this with \emph{neither … nor}, which is why the English needs two words too. A one-word translation would be a lie about the shape.
\end{grammarlens}

\subsection*{Your turn}
Say these aloud:

\begin{itemize}
\raggedright
  \item \emph{chahā āṇi dūdh} — both offered
  \item \emph{nā chahā nā dūdh} — both refused
  \item where \textbf{\mr{ना}} sits, against where \textbf{\mr{आणि}} sits — \textbf{in front of each half, not between them}
\end{itemize}

\begin{tcolorbox}[breakable,colback=teal!4,colframe=teal!35!black,title={Before you move on}]
How many times must \textbf{\mr{ना}} be said? (\textbf{Twice}.) Where does it go, compared with \textbf{\mr{आणि}}? (\textbf{In front of each half, not between them}.)
\end{tcolorbox}
\section[… nā?]{… \mr{ना}? (… nā?) — \textquotedblleft{}isn't it?\textquotedblright{}}

\label{lesson:MR-C31-naa-tag}



\begin{quote}
You have just used \textbf{\mr{ना}} twice in a row to refuse two things. The same syllable does something completely different when it comes \textbf{once}, and at the \textbf{end}.
\end{quote}

\begin{grammarlens}[title={Grammar Lens: hanging a check on the end}]
\begin{quote}
\textbf{\mr{चहा} \mr{आहे} \mr{ना}?} — \emph{chahā āhe nā?} — \textbf{There's tea, isn't there?}
\end{quote}

Take any statement you can already make. Put \textbf{\mr{ना}} on the end. It is now a check: you are telling the listener what you believe and inviting them to confirm it. The answer comes back as \textbf{\mr{हो}} or \textbf{\mr{नाही}}.

This is worth more than its size. A Marathi conversation is full of these, and they are how a beginner keeps a turn going without having to build a question from scratch. One syllable turns every sentence in this book into something a listener has to answer.

Notice how far the negative \textbf{\mr{ना}} has travelled: doubled in front, it wipes two things out; single at the end, it politely asks for a \emph{yes}. The tag is the old negative used as an invitation — the same move English makes in \emph{…, isn't it?}, where a negative asks for agreement.
\end{grammarlens}

\subsection*{Your turn}
Say these aloud:

\begin{itemize}
\raggedright
  \item \emph{chahā āhe nā?} then answer yourself — \emph{ho}
  \item \emph{malā marāṭhī yete nā?} — I know Marathi, don't I?
  \item \emph{nā chahā nā dūdh} — and hear the same syllable doing the opposite job
\end{itemize}

\begin{tcolorbox}[breakable,colback=teal!4,colframe=teal!35!black,title={Before you move on}]
What does one \textbf{\mr{ना}} at the END of a sentence do? (\textbf{Asks the listener to confirm it}.) And two \textbf{\mr{ना}} in FRONT? (\textbf{Refuses both things}.)
\end{tcolorbox}
\section[dusrā]{\mr{दुसरा} (dusrā) — the other one}

\label{lesson:MR-C31-dusra}



\begin{quote}
Count to five out loud. You have owned \textbf{\mr{एक}} since the numbers chapter. It has been waiting for a partner.
\end{quote}

\subsection*{You'll want to know: \mr{दुसरा}}
\begin{quote}
\textbf{\mr{दुसरा}} — \emph{dusrā} — \textbf{the other, the second}
\end{quote}

\textbf{\mr{दुसरा}} is \textbf{\mr{दोन}} (\textquotedblleft{}two\textquotedblright{}) wearing an ordinal ending: literally \emph{two-th}. Marathi builds it the way English builds \emph{second} out of nothing you can see, and Marathi's is the more honest of the two — you can still read the \emph{two} inside it.

Like every \textbf{-\mr{आ}} word in this book it agrees: \textbf{\mr{दुसरा}} for a masculine thing, \textbf{\mr{दुसरी}} feminine, \textbf{\mr{दुसरं}} neuter.

\begin{grammarlens}[title={Grammar Lens: \mr{एक} … \mr{दुसरा} … — sharing a pair out}]
\begin{quote}
\textbf{\mr{एक} \mr{चहा} \mr{आणि} \mr{दुसरा} \mr{पाणी}.} — \emph{ek chahā āṇi dusrā pāṇī.} — \textbf{One is tea and the other is water.}
\end{quote}

Set \textbf{\mr{एक}} in front of the first thing and \textbf{\mr{दुसरा}} in front of the second, and the pair is now shared out rather than merely listed. \textbf{\mr{आणि}} would only stack them; this frame says \emph{this one, and that one}. It is the last of the joining shapes that needs a word at BOTH ends, and it is how Marathi distributes a choice a speaker has already made.
\end{grammarlens}

\subsection*{Your turn}
Say these aloud:

\begin{itemize}
\raggedright
  \item \emph{ek chahā āṇi dusrā pāṇī}
  \item the number hiding inside \emph{dusrā} — \emph{don}
  \item the three agreeing forms — \emph{dusrā}, \emph{dusrī}, \emph{dusraṁ}
\end{itemize}

\begin{tcolorbox}[breakable,colback=teal!4,colframe=teal!35!black,title={Before you move on}]
What number is \textbf{\mr{दुसरा}} built on? (\textbf{\mr{दोन}}.) What does putting \textbf{\mr{एक}} and \textbf{\mr{दुसरा}} at the two ends do that \textbf{\mr{आणि}} alone cannot? (\textbf{Shares the pair out — this one, that one}.)

Sources: \href{https://en.wiktionary.org/wiki/\%E0\%A4\%A6\%E0\%A5\%81\%E0\%A4\%B8\%E0\%A4\%B0\%E0\%A4\%BE}{Wiktionary: \mr{दुसरा}}.
\end{tcolorbox}
\section[nā … nā · … nā? · ek … dusrā …]{Words that need two places}

\label{lesson:MR-R31-pairs}



\begin{quote}
Chapter thirty gave you three joiners that each need ONE place. Name them. Chapter thirty-one gave you three that need more.
\end{quote}

\subsection*{Your turn}
Say these aloud:

\begin{itemize}
\raggedright
  \item \emph{nā chahā nā dūdh}
  \item \emph{ek chahā āṇi dusrā pāṇī}
  \item \emph{chahā āhe nā?} and answer \emph{ho}
\end{itemize}

\begin{tcolorbox}[breakable,colback=teal!4,colframe=teal!35!black,title={Before you move on}]
Which of this chapter's three shapes puts its word at the END? (\textbf{The tag \mr{ना}}.) Which two put a word in FRONT of each half? (\textbf{\mr{ना} … \mr{ना}} and \textbf{\mr{एक} … \mr{दुसरा}}.) Then stop.
\end{tcolorbox}
